\uuid{Zm7C}
\exo7id{3625}
\titre{exo7 3625}
\auteur{quercia}
\organisation{exo7}
\datecreate{2010-03-10}
\isIndication{false}
\isCorrection{false}
\chapitre{Endomorphisme particulier}
\sousChapitre{Autre}
\cpge{mpsi-mp2i.espaces-vectoriels.c, pcsi.espaces-vectoriels.c, ptsi.espaces-vectoriels.c, mpsi-mp2i.dimension-finie.b, pcsi.dimension-finie.b, ptsi.dimension-finie.b, mpsi-mp2i.matrices-applications-lineaires.a, pcsi.matrices-applications-lineaires.a, ptsi.matrices-applications-lineaires.a, bcpst1.espace-vectoriel-kn.a, bcpst1.espace-vectoriel-kn.b}
\module{Algèbre}
\niveau{L2}
\difficulte{}

\contenu{
\texte{
Dans $ K^3$ on considère les formes linéaires :
$f_1(\vec x) = x+y-z$,
$f_2(\vec x) = x-y+z$,
$f_3(\vec x) = x+y+z$.
}
\begin{enumerate}
    \item \question{Montrer que $(f_1,f_2,f_3)$ est une base de $( K^3)^*$.}
    \item \question{Trouver la base duale.}
\end{enumerate}
}
